%% kapa: exam=1 round=28 year=2017 period=2 form=A subject=appraisal_law subject_ko=감정평가관계법규 seq=34 number=34
%% source: Q-Net 감정평가사 기출문제 (cst00309) artlSeq=5199603 file=extracted/1008/5199603/2017년 제28회 감정평가사 1차시험 문제지/2교시 A형.pdf
%% conversion: kordoc 4.24.51 -> tools/gen1.py (auto); review=proofread-sonnet; flags=none
\begin{question}{34}
부동산등기법령상 용익권 및 담보권에 관한 등기에 대한 설명으로 옳은 것은?
\begin{choices}
  \item 등기관이 지상권설정의 등기를 할 때 지상권의 범위는 등기원인에 그 약정이 있는 경우에만 기록한다.
  \item 등기관이 근저당권설정의 등기를 할 때 채권의 최고액은 등기원인에 그 약정이 있는 경우에만 기록한다.
  \item 등기관이 전세권설정의 등기를 할 때 위약금 또는 배상금은 등기원인에 그 약정이 있는 경우에만 기록한다.
  \item 등기관이 전세금반환채권의 일부 양도를 원인으로 한 전세권 일부이전등기를 할 때 양도액은 기록하지 않는다.
  \item 등기관이 동일한 채권에 관하여 여러 개의 부동산에 관한 권리를 목적으로 하는 저당권설정의 등기를 할 경우, 부동산이 3개 이상일 때에는 공동담보목록을 작성하여야 한다.
\end{choices}
\end{question}
\answer{3}
